% answers_short.tex —— 2014 年·四川卷·理科 · 简短答案速查
% 由 tools/make_short_answers.py 自 answers.tex 的 \gksolhead/\gkanswer 提取，
% 并与 2014/2014_数学真题_四川理_答案要点.md 逐题比对后生成，勿手改（改源 MD 后重跑）。

\begin{gknotice}
\textbf{使用说明}：本节只给最终答案便于快速核对，完整推导见解析版。答案为编写者独立推导，\textbf{非官方标准答案}。
\end{gknotice}

\gkpart{选择题}
\begin{center}\begin{tabular}{ccccccccccc}
\toprule
题号 & 1 & 2 & 3 & 4 & 5 & 6 & 7 & 8 & 9 & 10 \\
\midrule
答案 & A & C & A & D & C & B & D & B & A & B \\
\bottomrule
\end{tabular}\end{center}

\gkpart{填空题}
\begin{itemize}[leftmargin=2.6em]
  \item[\textbf{11.}] $-2\mathrm i$。分子分母同乘 $1-\mathrm i$，得 $\dfrac{(2-2\mathrm i)(1-\mathrm i)}2=\dfrac{-4\mathrm i}2=-2\mathrm i$。
  \item[\textbf{12.}] $1$。周期为 $2$，故 $f\left(\dfrac32\right)=f\left(-\dfrac12\right)=-4\times\dfrac14+2=1$。
  \item[\textbf{13.}] $60$。设垂足为 $D$，$BD=46\cot67^\circ$，$CD=46\cot30^\circ=46\sqrt3$，故 $BC=46(\sqrt3-\cot67^\circ)\approx 46(1.73-0.424)\approx 60.08$，精确到个位为 $60\ \mathrm m$。
  \item[\textbf{14.}] $5$。$A(0,0)$，$B(1,3)$，两直线恒垂直，$P$ 在以 $AB$ 为直径的圆上，$|PA|^2+|PB|^2=10$，故 $|PA|\cdot|PB|\leqslant 5$，$|PA|=|PB|=\sqrt5$ 时取等号。
  \item[\textbf{15.}] ①③④。①为满射的等价表述；②错，如 $\arctan x$ 有界但无最值；③对，否则 $f=(f+g)-g$ 也有界；④对，$a\neq 0$ 时 $f$ 无上界，只能 $a=0$，此时 $f=\dfrac{x}{x^2+1}$ 有最大值和最小值。
\end{itemize}

\gkpart{解答题}
\begin{itemize}[leftmargin=2.6em]
  \item[\textbf{16.}] (1) $\left[-\dfrac\pi4+\dfrac{2k\pi}3,\ \dfrac\pi{12}+\dfrac{2k\pi}3\right]$（$k\in\mathbf Z$）。由 $-\dfrac\pi2+2k\pi\leqslant 3x+\dfrac\pi4\leqslant\dfrac\pi2+2k\pi$ 解出；(2) $-\sqrt2$ 或 $-\dfrac{\sqrt5}2$。记 $u=\cos\alpha-\sin\alpha$、$v=\cos\alpha+\sin\alpha$，条件化为 $v=\dfrac45u^2v$：$v=0$ 给 $\alpha=\dfrac{3\pi}4$、$u=-\sqrt2$；$u^2=\dfrac54$ 给 $u=-\dfrac{\sqrt5}2$（$\alpha\approx 97.2^\circ$ 或 $172.8^\circ$），两支均满足第二象限条件，原题未排除任一支
  \item[\textbf{17.}] (1) $P(X=-200)=\dfrac18$，$P(X=10)=\dfrac38$，$P(X=20)=\dfrac38$，$P(X=100)=\dfrac18$。三次击鼓中出现音乐的次数 $Y\sim B\left(3,\dfrac12\right)$，按次数对应得分；(2) $\dfrac{511}{512}$。每盘至少出现一次音乐的概率为 $1-\dfrac18=\dfrac78$，三盘相互独立，故所求为 $1-\left(\dfrac18\right)^3=\dfrac{511}{512}$；(3) 每盘得分的期望 $E(X)=\dfrac{-200+30+60+100}8=-\dfrac54<0$，故盘数越多平均每盘越接近 $-\dfrac54$ 分、总分呈下降趋势；根源是概率 $\dfrac18$ 的 $-200$ 分罚分对期望的贡献 $-25$ 压过了其余三种得分
  \item[\textbf{18.}] (1) 证明。由三视图得底面 $BCD$ 为边长 $2$ 的正三角形，$A$ 在底面的射影是 $BD$ 的中点 $O$，$AO=\sqrt3$；以 $O$ 为原点建系，由 $\overrightarrow{MN}\cdot\overrightarrow{NP}=t-\dfrac12=0$（$P$ 分 $\overrightarrow{BC}$ 的比为 $t$）得 $t=\dfrac12$，故 $P$ 为 $BC$ 中点；(2) $\dfrac{\sqrt{10}}5$。两半平面内垂直于棱 $NP$ 的向量取 $\left(\dfrac{\sqrt3}4,\dfrac12,\dfrac{\sqrt3}4\right)$ 与 $(0,1,0)$，夹角余弦为 $\dfrac{2}{\sqrt{10}}=\dfrac{\sqrt{10}}5$
  \item[\textbf{19.}] (1) $S_n=n^2-3n$。由 $4b_7=2^{a_8}$ 得 $a_8=a_7+2$，故 $d=2$，$a_n=2n-4$；(2) $T_n=2-\dfrac{n+2}{2^n}$。切线截距 $a_2-\dfrac1{\ln2}=2-\dfrac1{\ln2}$ 给出 $a_2=2$，故 $a_n=n$，$b_n=2^n$，$\dfrac{a_n}{b_n}=\dfrac n{2^n}$，错位相减即得
  \item[\textbf{20.}] (1) $\dfrac{x^2}6+\dfrac{y^2}2=1$。$c=2$，由正三角形条件得 $a^2=3b^2$，结合 $a^2-b^2=4$ 得 $b^2=2$，$a^2=6$；(2) ① 证明；② $T(-3,1)$ 或 $T(-3,-1)$。设 $T(-3,m)$，垂线为 $x=-2+ms$，$y=s$，代入椭圆得 $(m^2+3)s^2-4ms-2=0$；$PQ$ 中点 $\dfrac{2}{m^2+3}(-3,m)$ 在直线 $OT$ 上；又 $\dfrac{|TF|}{|PQ|}=\dfrac{m^2+3}{2\sqrt6\sqrt{1+m^2}}\geqslant\dfrac1{\sqrt3}$，等号当 $m^2=1$
  \item[\textbf{21.}] (1) 当 $a\leqslant\dfrac12$ 时 $g_{\min}=1-b$；当 $\dfrac12<a<\dfrac{\mathrm e}2$ 时 $g_{\min}=2a-2a\ln(2a)-b$；当 $a\geqslant\dfrac{\mathrm e}2$ 时 $g_{\min}=\mathrm e-2a-b$。由 $g'(x)=\mathrm e^x-2a$ 的零点位置分三种情形求得；(2) $\mathrm e-2<a<1$。由 $f(1)=0$ 得 $b=\mathrm e-1-a$，且 $f(0)=0$；对 $[0,x_0]$ 与 $[x_0,1]$ 用罗尔定理知 $g=f'$ 在 $(0,1)$ 内有两个零点，而 $g''=\mathrm e^x>0$ 使 $g$ 严格凸，故 $g(0)>0$、$g(1)>0$，即 $a+2-\mathrm e>0$、$1-a>0$；反之 $g(0)>0$、$g(1)>0$ 时由介值定理确有零点，故取值范围为 $(\mathrm e-2,1)$
\end{itemize}
